\thispagestyle{plain}
\begin{fullwidth}
\begin{center}
    \LARGE
    \textcolor{itesodblue}{\textbf{\mytitle}}\\
    \Large
    \vspace{0.4cm}
    \textcolor{itesomblue}{\textbf{\myauthor}}\\
    \vspace{0.9cm}
    \textcolor{itesolblue}{\textbf{Resumen}}
\end{center}
Los datos estructurados como grafos describen un dominio como un conjunto de
entidades junto con las relaciones entre ellas, una forma que las
representaciones tabulares, en rejilla y secuenciales no capturan: los grafos
varían en tamaño, no admiten un ordenamiento canónico de sus elementos y
conectan cada entidad con un número distinto de otras. Una red neuronal
estándar, cuya dimensión de entrada y ordenamiento de coordenadas son fijos, no
puede operar sobre estos datos sin descartar la estructura relacional que porta
buena parte de la señal. Las redes neuronales de grafos (GNN) son arquitecturas
de redes neuronales diseñadas para operar directamente sobre grafos, respetando
estas propiedades en lugar de descartarlas.

Este trabajo constituye una referencia autocontenida sobre las GNN para grafos
homogéneos bajo el marco de paso de mensajes. Primero se establecen los
fundamentos: cómo los grafos representan datos relacionales, por qué el
perceptrón multicapa (MLP) fracasa ante ellos, y la invarianza y equivarianza a
permutaciones que debe satisfacer toda capa definida sobre un grafo. A
continuación, se desarrolla el marco de paso de mensajes como la formulación
unificadora, descompuesto en las funciones de mensaje, agregación y
actualización, junto con el campo receptivo que inducen capas sucesivas, la
distinción entre los escenarios transductivo e inductivo, y el techo de la
capacidad expresiva.

Se derivan tres arquitecturas representativas como instancias del marco, una por
cada nivel estándar de aprendizaje sobre grafos: la Graph Convolutional Network (GCN)
para clasificación de nodos, el Graph Auto-Encoder (GAE) para predicción de
enlaces, y la Graph Isomorphism Network (GIN) para clasificación de grafos. Cada
una se implementa en PyTorch Geometric sobre un conjunto de datos de referencia
estándar --- Cora, Amazon Computers y NCI1, respectivamente --- y se evalúa
frente a un modelo de referencia MLP que predice únicamente a partir de las
características de los nodos; en todas las tareas la arquitectura que usa el
grafo mejora a la que no lo usa. Por último, se examinan las limitaciones del
marco: sobresuavizado (\textit{oversmoothing}), sobrecompresión
(\textit{oversquashing}), heterofilia y la cota sobre la capacidad expresiva.
\end{fullwidth}